\documentclass[10pt,a4paper]{amsart}
\usepackage{amsmath,amssymb,amsthm,graphicx}
% For commutative diagrams we use Xy-pic
\usepackage[all]{xy}

% We use 2cell for 2-commutative diagrams.
\xyoption{2cell}
% Theorem environments.
%
\theoremstyle{plain}
\newtheorem{theorem}[subsection]{Theorem}
\newtheorem{proposition}[subsection]{Proposition}
\newtheorem{lemma}[subsection]{Lemma}

\theoremstyle{definition}
\newtheorem{definition}[subsection]{Definition}
\newtheorem{example}[subsection]{Example}
\newtheorem{exercise}[subsection]{Exercise}
\newtheorem{situation}[subsection]{Situation}

\theoremstyle{remark}
\newtheorem{remark}[subsection]{Remark}
\newtheorem{remarks}[subsection]{Remarks}

\numberwithin{equation}{subsection}

% Macros
%
\def\lim{\mathop{\mathrm{lim}}\nolimits}
\def\colim{\mathop{\mathrm{colim}}\nolimits}
\def\Spec{\mathop{\mathrm{Spec}}}
\def\Hom{\mathop{\mathrm{Hom}}\nolimits}
\def\Ext{\mathop{\mathrm{Ext}}\nolimits}
\def\SheafHom{\mathop{\mathcal{H}\!\mathit{om}}\nolimits}
\def\SheafExt{\mathop{\mathcal{E}\!\mathit{xt}}\nolimits}
\def\Sch{\mathit{Sch}}
\def\Mor{\mathop{\mathrm{Mor}}\nolimits}
\def\Ob{\mathop{\mathrm{Ob}}\nolimits}
\def\Sh{\mathop{\mathit{Sh}}\nolimits}
\def\NL{\mathop{N\!L}\nolimits}
\def\CH{\mathop{\mathrm{CH}}\nolimits}
\def\proetale{{pro\text{-}\acute{e}tale}}
\def\etale{{\acute{e}tale}}
\def\QCoh{\mathit{QCoh}}
\def\Ker{\mathop{\mathrm{Ker}}}
\def\Im{\mathop{\mathrm{Im}}}
\def\Coker{\mathop{\mathrm{Coker}}}
\def\Coim{\mathop{\mathrm{Coim}}}

% Boxtimes
%
\DeclareMathSymbol{\boxtimes}{\mathbin}{AMSa}{"02}

%
% Macros for moduli stacks/spaces
%
\def\QCohstack{\mathcal{QC}\!\mathit{oh}}
\def\Cohstack{\mathcal{C}\!\mathit{oh}}
\def\Spacesstack{\mathcal{S}\!\mathit{paces}}
\def\Quotfunctor{\mathrm{Quot}}
\def\Hilbfunctor{\mathrm{Hilb}}
\def\Curvesstack{\mathcal{C}\!\mathit{urves}}
\def\Polarizedstack{\mathcal{P}\!\mathit{olarized}}
\def\Complexesstack{\mathcal{C}\!\mathit{omplexes}}
% \Pic is the operator that assigns to X its picard group, usage \Pic(X)
% \Picardstack_{X/B} denotes the Picard stack of X over B
% \Picardfunctor_{X/B} denotes the Picard functor of X over B
\def\Pic{\mathop{\mathrm{Pic}}\nolimits}
\def\Picardstack{\mathcal{P}\!\mathit{ic}}
\def\Picardfunctor{\mathrm{Pic}}
\def\Deformationcategory{\mathcal{D}\!\mathit{ef}}
\newenvironment{reference}{}{}
\newenvironment{slogan}{}{}
\newenvironment{history}{}{}
\usepackage{cleveref}
\crefname{originaldefinition}{original definition}{original definitions}
\Crefname{originaldefinition}{Original definition}{Original definitions}
\makeatletter
\newlabel{definition-big-small-h}{{0EU1}{source}{}{sourcepointer-1}{}}
\newlabel{definition-big-small-h@cref}{{[originaldefinition][0][]0EU1}{[1][0][]source}{}{}{}}
\makeatother
\crefname{originallemma}{original lemma}{original lemmas}
\Crefname{originallemma}{Original lemma}{Original lemmas}
\makeatletter
\newlabel{sheaves-lemma-directed-colimits-sections}{{009F}{source}{}{sourcepointer-2}{}}
\newlabel{sheaves-lemma-directed-colimits-sections@cref}{{[originallemma][0][]009F}{[1][0][]source}{}{}{}}
\makeatother
\crefname{originallemma}{original lemma}{original lemmas}
\Crefname{originallemma}{Original lemma}{Original lemmas}
\makeatletter
\newlabel{lemma-refine-check-h}{{0EVI}{source}{}{sourcepointer-3}{}}
\newlabel{lemma-refine-check-h@cref}{{[originallemma][0][]0EVI}{[1][0][]source}{}{}{}}
\makeatother
\crefname{originallemma}{original lemma}{original lemmas}
\Crefname{originallemma}{Original lemma}{Original lemmas}
\makeatletter
\newlabel{descent-lemma-sheaf-condition-holds}{{03DT}{source}{}{sourcepointer-4}{}}
\newlabel{descent-lemma-sheaf-condition-holds@cref}{{[originallemma][0][]03DT}{[1][0][]source}{}{}{}}
\makeatother
\crefname{originalsection}{original section}{original sections}
\Crefname{originalsection}{Original section}{Original sections}
\makeatletter
\newlabel{example-one-generator}{{0EVG}{source}{}{sourcepointer-5}{}}
\newlabel{example-one-generator@cref}{{[originalsection][0][]0EVG}{[1][0][]source}{}{}{}}
\makeatother
\crefname{originallemma}{original lemma}{original lemmas}
\Crefname{originallemma}{Original lemma}{Original lemmas}
\makeatletter
\newlabel{lemma-funny-blow-up}{{0EVJ}{source}{}{sourcepointer-6}{}}
\newlabel{lemma-funny-blow-up@cref}{{[originallemma][0][]0EVJ}{[1][0][]source}{}{}{}}
\makeatother
\crefname{originallemma}{original lemma}{original lemmas}
\Crefname{originallemma}{Original lemma}{Original lemmas}
\makeatletter
\newlabel{limits-lemma-descend-invertible-modules}{{0B8W}{source}{}{sourcepointer-7}{}}
\newlabel{limits-lemma-descend-invertible-modules@cref}{{[originallemma][0][]0B8W}{[1][0][]source}{}{}{}}
\makeatother
\crefname{originallemma}{original lemma}{original lemmas}
\Crefname{originallemma}{Original lemma}{Original lemmas}
\makeatletter
\newlabel{more-morphisms-lemma-globally-generated-vanishing}{{0E7H}{source}{}{sourcepointer-8}{}}
\newlabel{more-morphisms-lemma-globally-generated-vanishing@cref}{{[originallemma][0][]0E7H}{[1][0][]source}{}{}{}}
\makeatother
\crefname{originaltheorem}{original theorem}{original theorems}
\Crefname{originaltheorem}{Original theorem}{Original theorems}
\makeatletter
\newlabel{coherent-theorem-formal-functions}{{02OC}{source}{}{sourcepointer-9}{}}
\newlabel{coherent-theorem-formal-functions@cref}{{[originaltheorem][0][]02OC}{[1][0][]source}{}{}{}}
\makeatother
\title{Vector-bundle descent across a two-generator almost-blowup square after Frobenius localization}
\author{The Stacks Project Authors}
\date{stacks-project:flat.tex:local-snapshot-20261011; GNU FDL 1.2 or later}
\begin{document}
\maketitle
\noindent\textit{Editorial note.} The square is exactly Exampletwo-generators, not an arbitrary proper square. The originalcolim\_F category convention is included.\par\medskip
\noindent\textit{Editorial note.} Public736 concerns pseudo-coherent derived triples across affine thickenings. This target concerns vector bundles across an almost blowup in characteristicp after Frobenius localization.\par\medskip
\noindent\textit{Allowed standard prerequisites (editorial).}\begin{itemize}
\item Explicit prerequisite boundary. Characteristic-p Frobenius pullback, finite locally free modules, the explicitly defined two-generator almost-blowup square, Noetherian approximation, projective-line cohomology and formal functions. Function-levelFrobenius descent and thecolimF gluing lemma are retained.
\end{itemize}\medskip
\ifdefined\hypertarget\hypertarget{sourcepointer-1}{}\fi
\noindent\textit{Source reference.} Original definition 0EU1: focused-seeds/AGP/sources/stacks-snapshot-20261011/tags.txt, source lines 18536--18536. Stable Stacks tag for an expressly allowed foundational or notational input; the counted construction remains the two concrete core obligations. The retained local core is not replaced by this pointer. See standard\_prerequisites.\par\medskip
\ifdefined\hypertarget\hypertarget{sourcepointer-2}{}\fi
\noindent\textit{Source reference.} Original lemma 009F: focused-seeds/AGP/sources/stacks-snapshot-20261011/tags.txt, source lines 351--351. Stable Stacks tag for an expressly allowed foundational or notational input; the counted construction remains the two concrete core obligations. The retained local core is not replaced by this pointer. See standard\_prerequisites.\par\medskip
\ifdefined\hypertarget\hypertarget{sourcepointer-3}{}\fi
\noindent\textit{Source reference.} Original lemma 0EVI: focused-seeds/AGP/sources/stacks-snapshot-20261011/tags.txt, source lines 18588--18588. Stable Stacks tag for an expressly allowed foundational or notational input; the counted construction remains the two concrete core obligations. The retained local core is not replaced by this pointer. See standard\_prerequisites.\par\medskip
\ifdefined\hypertarget\hypertarget{sourcepointer-4}{}\fi
\noindent\textit{Source reference.} Original lemma 03DT: focused-seeds/AGP/sources/stacks-snapshot-20261011/tags.txt, source lines 4432--4432. Stable Stacks tag for an expressly allowed foundational or notational input; the counted construction remains the two concrete core obligations. The retained local core is not replaced by this pointer. See standard\_prerequisites.\par\medskip
\ifdefined\hypertarget\hypertarget{sourcepointer-5}{}\fi
\noindent\textit{Source reference.} Original section 0EVG: focused-seeds/AGP/sources/stacks-snapshot-20261011/tags.txt, source lines 18586--18586. Stable Stacks tag for an expressly allowed foundational or notational input; the counted construction remains the two concrete core obligations. The retained local core is not replaced by this pointer. See standard\_prerequisites.\par\medskip
\ifdefined\hypertarget\hypertarget{sourcepointer-6}{}\fi
\noindent\textit{Source reference.} Original lemma 0EVJ: focused-seeds/AGP/sources/stacks-snapshot-20261011/tags.txt, source lines 18589--18589. Stable Stacks tag for an expressly allowed foundational or notational input; the counted construction remains the two concrete core obligations. The retained local core is not replaced by this pointer. See standard\_prerequisites.\par\medskip
\ifdefined\hypertarget\hypertarget{sourcepointer-7}{}\fi
\noindent\textit{Source reference.} Original lemma 0B8W: focused-seeds/AGP/sources/stacks-snapshot-20261011/tags.txt, source lines 14154--14154. Stable Stacks tag for an expressly allowed foundational or notational input; the counted construction remains the two concrete core obligations. The retained local core is not replaced by this pointer. See standard\_prerequisites.\par\medskip
\ifdefined\hypertarget\hypertarget{sourcepointer-8}{}\fi
\noindent\textit{Source reference.} Original lemma 0E7H: focused-seeds/AGP/sources/stacks-snapshot-20261011/tags.txt, source lines 17782--17782. Stable Stacks tag for an expressly allowed foundational or notational input; the counted construction remains the two concrete core obligations. The retained local core is not replaced by this pointer. See standard\_prerequisites.\par\medskip
\ifdefined\hypertarget\hypertarget{sourcepointer-9}{}\fi
\noindent\textit{Source reference.} Original theorem 02OC: focused-seeds/AGP/sources/stacks-snapshot-20261011/tags.txt, source lines 3511--3511. Stable Stacks tag for an expressly allowed foundational or notational input; the counted construction remains the two concrete core obligations. The retained local core is not replaced by this pointer. See standard\_prerequisites.\par\medskip
\begin{lemma}
\label{lemma-vector-bundle-II}
Let $p$ be a prime number. Consider an almost blowup square $X, X', Z, E$
in characteristic $p$ as in Example \ref{example-two-generators}.
Then the functor
$$
G :
\colim_F \textit{Vect}(X)
\longrightarrow
\colim_F \textit{Vect}(Z)
\times_{\colim_F \textit{Vect}(E)}
\colim_F \textit{Vect}(X')
$$
is an equivalence.
\end{lemma}

\begin{example}
\label{example-two-generators}
Let $A$ be a ring. Let $f_1, f_2 \in A$ be elements.
$$
\xymatrix{
E = \text{Proj}(A/(f_1, f_2)[T_0, T_1]) \ar[r] \ar[d] &
\text{Proj}(A[T_0, T_1]/(f_2 T_0 - f_1 T_1) = X' \ar[d] \\
Z = \Spec(A/(f_1, f_2)) \ar[r] & \Spec(A) = X
}
$$
is an almost blowup square.
\end{example}

\section{Descent vector bundles in positive characteristic}
\label{section-descent-vectorbundles-h}

\noindent
A reference for this section is \cite{Witt-Grass}.

\medskip\noindent
For a scheme $S$ let us denote $\textit{Vect}(S)$ the category
of finite locally free $\mathcal{O}_S$-modules.
Let $p$ be a prime number. Let $S$ be a quasi-compact and quasi-separated
scheme over $\mathbf{F}_p$. In this section we will work with the
category
$$
\colim_F \textit{Vect}(S) =
\colim
\left(
\textit{Vect}(S) \xrightarrow{F^*}
\textit{Vect}(S) \xrightarrow{F^*}
\textit{Vect}(S) \xrightarrow{F^*} \ldots
\right)
$$
where $F : S \to S$ is the absolute Frobenius morphism. In down
to earth terms an object of this category is a pair $(\mathcal{E}, n)$
where $\mathcal{E}$ is a finite locally free $\mathcal{O}_S$-module
and $n \geq 0$ is an integer. For morphisms we take
$$
\Hom_{\colim_F \textit{Vect}(S)}((\mathcal{E}, n), (\mathcal{G}, m)) =
\colim_N \Hom_S(F^{N - n, *}\mathcal{E}, F^{N - m, *}\mathcal{G})
$$
where $F : S \to S$ is the absolute Frobenius morphism of $S$.
Thus the object $(\mathcal{E}, n)$ is isomorphic to the object
$(F^*\mathcal{E}, n + 1)$.

\begin{lemma}
\label{lemma-h-sheaf-colim-F}
Let $p$ be a prime number. Let $S$ be a scheme over $\mathbf{F}_p$.
Let $(\Sch/S)_h$ be a site as in Definition \ref{definition-big-small-h}.
There is a unique sheaf $\mathcal{F}$ on $(\Sch/S)_h$ such that
$$
\mathcal{F}(X) = \colim_F \Gamma(X, \mathcal{O}_X)
$$
for any quasi-compact and quasi-separated object $X$ of $(\Sch/S)_h$.
\end{lemma}

\begin{proof}
Denote $\mathcal{F}$ the Zariski sheafification
of the functor
$$
X \longrightarrow \colim_F \Gamma(X, \mathcal{O}_X)
$$
For quasi-compact and quasi-separated schemes $X$
we have $\mathcal{F}(X) = \colim_F \Gamma(X, \mathcal{O}_X)$.
by Sheaves, Lemma \ref{sheaves-lemma-directed-colimits-sections}
and the fact that $\mathcal{O}$ is a sheaf for the Zariski topology.
Thus it suffices to show that $\mathcal{F}$ is a h sheaf.
To prove this we check conditions (1), (2), (3), and (4) of
Lemma \ref{lemma-refine-check-h}.
Condition (1) holds because we performed an (almost unnecessary)
Zariski sheafification. Condition (2) holds because
$\mathcal{O}$ is an fppf sheaf (Descent, Lemma
\ref{descent-lemma-sheaf-condition-holds}) and if
$A$ is the equalizer of two maps $B \to C$ of $\mathbf{F}_p$-algebras,
then $\colim_F A$ is the equalizer of the two maps
$\colim_F B \to \colim_F C$.

\medskip\noindent
We check condition (3). Let $A, f, J$ be as in
Example \ref{example-one-generator}.
We have to show that
$$
\colim_F A = \colim_F A/J \times_{\colim_F A/fA + J} \colim_F A/fA
$$
This reduces to the following algebra question: suppose $a', a'' \in A$
are such that $F^n(a' - a'') \in fA + J$. Find $a \in A$ and $m \geq 0$
such that $a - F^m(a') \in J$ and $a - F^m(a'') \in fA$ and show that
the pair $(a, m)$ is uniquely determined up to a replacement of the
form $(a, m) \mapsto (F(a), m + 1)$.
To do this just write $F^n(a' - a'') = f h + g$ with $h \in A$ and $g \in J$
and set $a = F^n(a') - g = F^n(a'') + fh$ and set $m = n$.
To see uniqueness, suppose $(a_1, m_1)$ is a second solution.
By a replacement of the form given above we may assume $m = m_1$.
Then we see that $a - a_1 \in J$ and $a - a_1 \in fA$.
Since $J$ is annihilated by a power of $f$ we see that
$a - a_1$ is a nilpotent element. Hence $F^k(a - a_1)$ is zero
for some large $k$. Thus after doing more replacements we get
$a = a_1$.

\medskip\noindent
We check condition (4). Let $X, X', Z, E$ be as in
Example \ref{example-two-generators}. By
Lemma \ref{lemma-funny-blow-up} we see that
$$
\mathcal{F}(X) = \colim_F \Gamma(X, \mathcal{O}_X)
\longrightarrow
\colim_F \Gamma(X', \mathcal{O}_{X'}) = \mathcal{F}(X')
$$
is bijective. Since $E = \mathbf{P}^1_Z$ in this case we also
see that $\mathcal{F}(Z) \to \mathcal{F}(E)$ is bijective.
Thus the conclusion holds in this case as well.
\end{proof}

\begin{lemma}
\label{lemma-colim-F-Vect}
Let $p$ be a prime number. Let $S$ be a quasi-compact and quasi-separated
scheme over $\mathbf{F}_p$. The category $\colim_F \textit{Vect}(S)$
is equivalent to the category of finite locally free modules
over the sheaf of rings $\colim_F \mathcal{O}_S$ on $S$.
\end{lemma}

\begin{proof}
Omitted.
\end{proof}

\begin{proof}
Fully faithfulness. Suppose that $(\mathcal{E}, n)$ and
$(\mathcal{F}, m)$ are objects of $\colim_F \textit{Vect}(X)$.
Let $(a, b) : G(\mathcal{E}, n) \to G(\mathcal{F}, m)$
be a morphism in the RHS. We may choose $N \gg 0$ and
think of $a$ as a map
$a : F^{N - n, *}\mathcal{E}|_Z \to F^{N - m, *}\mathcal{F}|_Z$
and $b$ as a map
$b : F^{N - n, *}\mathcal{E}|_{X'} \to F^{N - m, *}\mathcal{F}|_{X'}$
agreeing over $E$.
Choose a finite affine open covering
$X = X_1 \cup \ldots \cup X_n$ such that $\mathcal{E}|_{X_i}$
and $\mathcal{F}|_{X_i}$ are finite free $\mathcal{O}_{X_i}$-modules.
For each $i$ the base change
$$
\xymatrix{
E_i \ar[r] \ar[d] & X'_i \ar[d] \\
Z_i \ar[r] & X_i
}
$$
is another almost blow up square as in Example \ref{example-two-generators}.
For these squares we know that
$$
\colim_F H^0(X_i, \mathcal{O}_{X_i}) =
\colim_F H^0(Z_i, \mathcal{O}_{Z_i})
\times_{\colim_F H^0(E_i, \mathcal{O}_{E_i})}
\colim_F H^0(X'_i, \mathcal{O}_{X'_i})
$$
by Lemma \ref{lemma-h-sheaf-colim-F} (see proof of the lemma).
Hence after increasing $N$ we may assume
the maps $a|_{Z_i}$ and $b|_{X'_i}$ come from
maps $c_i : F^{N - n, *}\mathcal{E}|_{X_i} \to F^{N - m, *}\mathcal{F}|_{X_i}$.
After possibly increasing $N$ we may assume $c_i$ and $c_j$
agree on $X_i \cap X_j$. Thus these maps glue to give the
desired morphism $(\mathcal{E}, n) \to (\mathcal{F}, m)$
in the LHS.

\medskip\noindent
Essential surjectivity. Let $(\mathcal{F}, \mathcal{G}, \varphi)$ be a
triple consisting of
a finite locally free $\mathcal{O}_Z$-module $\mathcal{F}$,
a finite locally free $\mathcal{O}_{X'}$-module $\mathcal{G}$, and
an isomorphism $\varphi : \mathcal{F}|_E \to \mathcal{G}|_E$.
We have to show that after replacing this triple by a Frobenius
power pullback, it comes from a finite locally free $\mathcal{O}_X$-module.

\medskip\noindent
Noetherian reduction; we urge the reader to skip this paragraph.
Recall that $X = \Spec(A)$ and $Z = \Spec(A/(f_1, f_2))$,
$X' = \text{Proj}(A[T_0, T_1]/(f_2T_0 - f_1T_1))$, and
$E = \mathbf{P}^1_Z$. By Limits, Lemma
\ref{limits-lemma-descend-invertible-modules}
we can find a finitely generated $\mathbf{F}_p$-subalgebra
$A_0 \subset A$ containing $f_1$ and $f_2$ such that the triple
$(\mathcal{F}, \mathcal{G}, \varphi)$ descends to
$X_0 = \Spec(A_0)$ and $Z_0 = \Spec(A_0/(f_1, f_2))$,
$X_0' = \text{Proj}(A_0[T_0, T_1]/(f_2T_0 - f_1T_1))$, and
$E_0 = \mathbf{P}^1_{Z_0}$. Thus we may assume our schemes
are Noetherian.

\medskip\noindent
Assume $X$ is Noetherian. We may choose a finite affine open covering
$X = X_1 \cup \ldots \cup X_n$ such that $\mathcal{F}|_{Z \cap X_i}$ is free.
Since we can glue objects of $\colim_F \textit{Vect}(X)$
in the Zariski topology (Lemma \ref{lemma-colim-F-Vect}), and
since we already know
fully faithfulness over $X_i$ and $X_i \cap X_j$ (see first paragraph
of the proof), it suffices to prove the existence over each $X_i$.
This reduces us to the case discussed in the next paragraph.

\medskip\noindent
Assume $X$ is Noetherian and $\mathcal{F} = \mathcal{O}_Z^{\oplus r}$.
Using $\varphi$ we get an isomorphism
$\mathcal{O}_E^{\oplus r} \to \mathcal{G}|_E$.
Let $I = (f_1, f_2) \subset A$.
Let $\mathcal{I} \subset \mathcal{O}_{X'}$
be the ideal sheaf of $E$; it is globally generated by $f_1$ and $f_2$.
For any $n$ there is a surjection
$$
(\mathcal{I}^n/\mathcal{I}^{n + 1})^{\oplus r} =
\mathcal{I}^n/\mathcal{I}^{n + 1} \otimes_{\mathcal{O}_E}
\mathcal{G}|_E \longrightarrow
\mathcal{I}^n\mathcal{G}/\mathcal{I}^{n + 1}\mathcal{G}
$$
Hence the first cohomology group of this module is zero.
Here we use that $E = \mathbf{P}^1_Z$ and hence its structure
sheaf and in fact any globally generated quasi-coherent module
has vanishing $H^1$. Compare with More on Morphisms, Lemma
\ref{more-morphisms-lemma-globally-generated-vanishing}.
Then using the short exact sequences
$$
0 \to  \mathcal{I}^n\mathcal{G}/\mathcal{I}^{n + 1}\mathcal{G} \to
\mathcal{G}/\mathcal{I}^{n + 1}\mathcal{G} \to
\mathcal{G}/\mathcal{I}^n\mathcal{G} \to 0
$$
and induction, we see that
$$
\lim H^0(X', \mathcal{G}/\mathcal{I}^n\mathcal{G})
\to
H^0(E, \mathcal{G}|_E) = H^0(E, \mathcal{O}_E^{\oplus r}) =
A/I^{\oplus r}
$$
is surjective. By the theorem on formal functions
(Cohomology of Schemes, Theorem \ref{coherent-theorem-formal-functions})
this implies that
$$
H^0(X', \mathcal{G}) \to
H^0(E, \mathcal{G}|_E) = H^0(E, \mathcal{O}_E^{\oplus r}) =
A/I^{\oplus r}
$$
is surjective. Thus we can choose a map
$\alpha : \mathcal{O}_{X'}^{\oplus r} \to \mathcal{G}$
which is compatible with the given trivialization
of $\mathcal{G}|_E$. Thus $\alpha$ is an isomorphism over
an open neighbourhood of $E$ in $X'$. Thus every point
of $Z$ has an affine open neighbourhood where
we can solve the problem. Since $X' \setminus E \to X \setminus Z$
is an isomorphism, the same holds for points of $X$ not in $Z$.
Thus another Zariski glueing argument finishes the proof.
\end{proof}

\bibliographystyle{plain}
\bibliography{my}
\end{document}
